\section{Does prediction help acting? The contrast gap}
\label{sec:gap}

The survey's question can be stated precisely. Let a \emph{direct policy} map observations to actions,
and let a \emph{world-model policy} first predict future observations and then plan an action against
that prediction. The advantage of prediction is the closed-loop task-success \emph{gain} of the
world-model policy over a matched direct policy, and it is only meaningful when reported per
capability, since prediction should help more where hidden state, occlusion, or long horizons make
foresight valuable. Figure~\ref{fig:eval_loop} draws the evaluation loop this definition requires: a
task forks into a direct-VLA branch and a world-model branch, both run in one closed loop with state
feedback, both scored by the same task-success measure, and the difference read off per capability.

% The evaluation loop that isolates the advantage of prediction (P1 analog: fig_architectures).
% Purely structural (no counts): shows where a benchmark must run BOTH a direct-VLA policy and a
% world-model policy under one closed loop, and where the per-capability advantage is measured.
\begin{figure}[t]
\centering
\resizebox{\textwidth}{!}{%
\begin{tikzpicture}[
 box/.style={draw=black!55, rounded corners, align=center, font=\footnotesize\bfseries,
             minimum height=9mm, text width=20mm, inner sep=3pt},
 io/.style={draw=none, font=\footnotesize\bfseries\itshape, align=center},
 gap/.style={box, fill=secEl, draw=secE!60!black, line width=1.1pt, text width=24mm},
 ar/.style={-{Stealth[length=2mm]}, darkgray, line width=0.8pt},
 fb/.style={-{Stealth[length=1.8mm]}, gray!70, line width=0.6pt},
]
 \node[io] (task) at (0,0) {task /\\ scene};
 \node[box, fill=cat-vla] (vla) at (3.1,1.15) {direct VLA\\ policy};
 \node[box, fill=secFl] (wm) at (3.1,-1.15) {world model:\\ predict $\rightarrow$ plan};
 \node[box, fill=secCl] (env) at (6.4,0) {environment\\ (closed loop)};
 \node[box, fill=secBl] (succ) at (9.4,0) {task\\ success};
 \node[gap] (adv) at (12.5,0) {advantage of\\ prediction ($\delta$)\\ per capability ($\beta$)};

 \draw[ar] (task) -- (vla);
 \draw[ar] (task) -- (wm);
 \draw[ar] (vla) -- (env);
 \draw[ar] (wm) -- (env);
 \draw[ar] (env) -- (succ);
 \draw[ar] (succ) -- (adv);
 % state-feedback loops (what makes it closed-loop); label once on top
 \draw[fb] (env.north) to[out=90,in=90] node[above,font=\scriptsize\bfseries,text=gray!20!black]{state feedback} (vla.north);
 \draw[fb] (env.south) to[out=-90,in=-90] (wm.south);
 % gap callout: the two ways today's suites break the loop
 \node[draw=BrickRed!55!black, dashed, rounded corners, fill=BrickRed!5, inner sep=4pt,
       text width=30mm, font=\scriptsize\bfseries, text=BrickRed!85!black, align=center] at (3.1,-3.1)
   {model-agnostic suites run \textbf{one} branch (no contrast); world-model suites never reach the loop};
 \draw[BrickRed!55!black, dashed, -{Stealth[length=1.6mm]}] (3.1,-2.2) -- (3.1,-2.55);
\end{tikzpicture}}
\caption{\textbf{The evaluation loop that isolates the advantage of prediction.} A benchmark that
can answer ``does world modelling help a robot act'' must run a direct Vision-Language-Action policy
\emph{and} a world-model policy (predict then plan) under one closed loop, then report the
closed-loop gain ($\delta$) sliced by capability ($\beta$). Today's suites break this loop in one of
two ways: model-agnostic policy suites host a single policy and never build the contrast, while
world-model suites score the prediction open-loop and never execute it (\S\ref{sec:gap}).}
\Description{A left-to-right diagram. A task/scene forks into two branches, a direct VLA policy and a
world model that predicts then plans; both feed an environment closed loop with state feedback,
producing task success, which feeds an advantage-of-prediction measurement sliced per capability. An
annotation notes that model-agnostic suites run only one branch and world-model suites never close
the loop, so the contrast is missing.}
\label{fig:eval_loop}
\end{figure}


Today's benchmarks break this loop in one of two ways, visible as the two poles of
Figure~\ref{fig:corpus_collage}. Model-agnostic policy and embodied suites (Sections~\ref{sec:lane_policy},
\ref{sec:lane_embodied}) close the loop but run a single branch, so there is no contrast to read.
World-model-evaluation suites (Section~\ref{sec:lane_wmeval}) score the prediction branch open-loop
and never execute it, so a high score certifies generation quality, not task success. The
prediction-to-action bridges (Section~\ref{sec:lane_bridge}) close the loop for a world-model policy
but do not yet run the matched VLA branch per capability. No lane completes the diagram.

\subsection{``World model'' is three things, none of them sufficient alone}

Part of why the contrast is hard to build is that the term ``world model'' is overloaded.
Table~\ref{tab:wm_senses} separates two axes the literature routinely conflates. Along a
\emph{representation} axis a world model may be a latent-dynamics predictor (LWM; DreamerV3, TD-MPC2),
an action-conditioned future-frame predictor (WAM; iVideoGPT, GR-2), or a generative video model
(WFM; Sora, Cosmos, Genie). Along a \emph{role} axis it may be a neural simulator that trains or
evaluates a policy, or the object that a benchmark scores. The crucial fact is that no single sense
is simultaneously action-conditioned, run in a closed loop, and readily comparable to a VLA policy
under one protocol: latent and action-conditioned models act but are rarely scored by a standard
generation metric, while generative models carry the metrics but are rarely closed into a loop. A
benchmark that contrasts a world-model policy with a VLA policy therefore has to assemble properties
that no existing sense of the term supplies at once.

% The senses in which the literature uses the term "world model" (P1 analog: Table 5, skill senses).
% Split along the two axes a parallel audit surfaced: representation (what it is) vs role (what it is for).
\begin{table*}[t]
\centering
\footnotesize
\setlength{\tabcolsep}{3.5pt}
\renewcommand{\arraystretch}{1.25}
\resizebox{\textwidth}{!}{%
\begin{tabular}{@{}p{0.16\textwidth} p{0.23\textwidth} p{0.18\textwidth} c c c c@{}}
\toprule
\textbf{Sense of ``world model''} & \textbf{What it is} & \textbf{Examples} &
\textbf{Act.\ cond.} & \textbf{Loop} & \textbf{Metric} & \textbf{VLA-cmp.} \\
\midrule
\multicolumn{7}{@{}l}{\rule{0pt}{2.4ex}\emph{As a representation (what the model is)}}\\
latent WM (LWM) & latent-dynamics predictor $z_{t+1}=f(z_t,a_t)$ & DreamerV3, TD-MPC2 & \yy & \yy & \pp & \pp \\
action-cond.\ WM (WAM) & action-conditioned future-frame predictor & iVideoGPT, GR-2 & \yy & \pp & \pp & \pp \\
generative WM (WFM) & text / image-to-video generator & Sora, Cosmos, Genie & \pp & \nn & \yy & \nn \\
\multicolumn{7}{@{}l}{\rule{0pt}{2.6ex}\emph{As a role (what the model is used for)}}\\
neural simulator & learned environment for policy training / eval & UniSim & \yy & \yy & \pp & \pp \\
evaluation target & benchmarks that \emph{score} world models & WorldSimBench, WorldModelBench, EWMBench & \dm & \pp & \yy & \nn \\
\bottomrule
\end{tabular}}
\caption{The term ``world model'' spans two orthogonal axes the literature routinely conflates: a
\emph{representation} axis (what the model is: latent-state dynamics, action-conditioned video, or
generative video, a soft boundary since action-conditioned generators straddle WAM and WFM) and a
\emph{role} axis (what it is used for: a neural simulator, or the object a benchmark scores).
Columns record whether a sense is \textbf{act}ion-\textbf{cond}itioned, is used in a closed
\textbf{loop}, carries a standard evaluation \textbf{metric}, and is readily \textbf{VLA-comp}arable
under one protocol (\yy\ yes, \pp\ partial, \nn\ no, \dm\ n/a). No single sense is action-conditioned,
closed-loop, and VLA-comparable at once, which is why a benchmark that contrasts a world-model policy
with a VLA policy per capability (\S\ref{sec:gap}) is still missing. Evaluation benchmarks score the
world model as an \emph{object}; they are not world models acting as judges.}
\label{tab:wm_senses}
\end{table*}


\subsection{Why each lane leaves the contrast unbuilt}

Table~\ref{tab:branch_limits} compares the four lanes on six evaluation axes and makes the pattern
lane-by-lane. Policy and embodied suites score the right quantity, closed-loop success, on the right
capabilities, but hold the model family fixed and so never contrast it. World-model-evaluation suites
vary the model family richly but score the wrong quantity for our question, open-loop prediction
quality, and never close the loop. The bridges are the only lane that scores closed-loop success of a
world model, but they are four benchmarks and none slices a VLA head-to-head by capability. The
result is a coverage hole precisely where the survey's question lives: the intersection of
closed-loop scoring, per-capability slicing, and a world-model-versus-VLA family contrast.

% Evaluation-lane comparison matrix: 4 lanes x 6 evaluation axes
% (loop closure, task horizon, capability breadth, model-family contrast, reproducibility, blind spot).
\begin{table*}[t]
\centering
\scriptsize
\setlength{\tabcolsep}{3pt}
\renewcommand{\arraystretch}{1.45}
\resizebox{\textwidth}{!}{%
\begin{tabular}{@{}>{\raggedright\arraybackslash}p{0.135\textwidth} >{\raggedright\arraybackslash}p{0.195\textwidth} >{\raggedright\arraybackslash}p{0.195\textwidth} >{\raggedright\arraybackslash}p{0.195\textwidth} >{\raggedright\arraybackslash}p{0.195\textwidth}@{}}
\toprule
 & \textbf{Policy suites}~\cite{libero} & \textbf{Embodied nav / social}~\cite{habitat30}
 & \textbf{World-model eval}~\cite{worldmodelbench} & \textbf{Bridges}~\cite{worldinworld} \\
\midrule
\textbf{Loop closure} &
\yy\ closed-loop; steps the environment~\cite{calvin} &
\yy\ closed-loop; long-horizon rollouts~\cite{behavior1k} &
\nn\ open-loop; scores generation, no action~\cite{physicsiq} &
\pp\ converts a prediction into an executed action~\cite{worldsimbench} \\
\textbf{Task horizon} &
\pp\ short reactive skills; long by chaining~\cite{libero} &
\yy\ long-horizon household and dialog~\cite{teach} &
\pp\ short clips; temporal but not task-level~\cite{ewmbench} &
\pp\ a single executed rollout~\cite{robowmbench} \\
\textbf{Capability breadth} &
\pp\ manipulation + transfer; hidden-state rare~\cite{thecolosseum} &
\yy\ navigation, social, some hidden-state~\cite{robothor} &
\pp\ physical / temporal; counterfactual barely~\cite{worldprediction} &
\nn\ narrow; executability only~\cite{worldarena} \\
\textbf{Model-family contrast} &
\nn\ model-agnostic host of any policy~\cite{roboarena} &
\pp\ memory-vs-Markovian in one suite~\cite{liberomem} &
\nn\ world-model-only scores~\cite{worldmodelbench} &
\pp\ WM as policy substrate; no VLA pole~\cite{worldinworld} \\
\textbf{Reproducibility} &
\pp\ sim; some real-vs-sim transfer~\cite{simpler} &
\pp\ sim-heavy; a sim-to-real gap~\cite{robothor} &
\yy\ fixed offline sets; deterministic scoring~\cite{vbench20} &
\nn\ new, still-unstandardised protocols~\cite{worldarena} \\
\textbf{Structural blind spot} &
hosts a world-model policy but never builds the VLA contrast &
capability cuts exist, but no per-capability family ablation &
visual quality is not task success; no behavioural utility &
reaches action, but omits the per-capability VLA head-to-head \\
\bottomrule
\end{tabular}}
\caption{\textbf{Evaluation-lane comparison matrix}: the four benchmark lanes on six evaluation
axes (loop closure, task horizon, capability breadth, model-family contrast, reproducibility, and
characteristic blind spot). Marks grade each lane on the axis (\yy\ favourable, \pp\ mixed, \nn\
weak), with an anchoring benchmark per cell. No lane closes the loop, slices capability, \emph{and}
builds a native VLA-vs-world-model contrast at once; the bridge lane comes closest yet still omits
the per-capability VLA head-to-head, which is the empty cell this survey names
(\S\ref{sec:gap}).}
\label{tab:branch_limits}
\end{table*}


\subsection{Measuring the advantage: assessed, not adopted}

Closing the loop is necessary but not sufficient; the gain also needs an instrument. Proposals exist,
including behavioural advantage-of-prediction scores such as ACTION-ATLAS's World Advantage Score,
but they are put forward in position papers~\cite{pos_yangyu} rather than established across a
benchmark suite. We therefore \emph{assess} such instruments as candidates rather than adopting any
one as ground truth, and Section~\ref{sec:future} specifies how an advantage-of-prediction quantity
should be reported: as a per-capability distribution or curve against a matched baseline, never as a
single headline number. The empirical scale of the gap is worth restating: of 160 benchmarks, 138
are model-agnostic, 11 build a partial contrast, and 11 build an explicit one; counterfactual
capability, where foresight should matter most, is almost entirely unmeasured; and only four
benchmarks reach executed action at all.

\FloatBarrier
